\section{\merit}
\label{sec:method}

The name states the goal: every input should earn its share of the result.
\merit{} has two branches that answer two different questions. The
\textbf{coupling branch} asks whether the recording follows a candidate's speech
over time. Only EEG enters it, because only EEG has a timed relationship with a
speech signal. The \textbf{orientation branch} asks which loudspeaker the
listener was facing. It receives no audio, so it cannot use audio cues.
Figure~\ref{fig:arch} shows the whole decoder.

\begin{figure}[t]
\centering
\definecolor{cIn}{HTML}{E8EFFB}\definecolor{cInB}{HTML}{3B6FB6}
\definecolor{cCo}{HTML}{E4EDFA}\definecolor{cCoB}{HTML}{3B6FB6}
\definecolor{cOr}{HTML}{E4F1EA}\definecolor{cOrB}{HTML}{3F8560}
\definecolor{cFz}{HTML}{F1EAF7}\definecolor{cFzB}{HTML}{8E5BA6}
\definecolor{cAd}{HTML}{FBEDE1}\definecolor{cAdB}{HTML}{BB6C24}
\resizebox{0.72\textwidth}{!}{%
\begin{tikzpicture}[font=\scriptsize,>=Stealth,x=1cm,y=1cm,
  nd/.style={draw,rounded corners=2pt,align=center,line width=0.6pt,inner sep=3pt},
  inp/.style={nd,fill=cIn,draw=cInB,minimum width=20mm,minimum height=5.6mm},
  enc/.style={nd,fill=white,draw=gray!60,minimum width=21mm,minimum height=5.6mm},
  frz/.style={nd,fill=cFz,draw=cFzB,minimum width=21mm,minimum height=5.6mm},
  fl/.style={->,line width=0.6pt,gray!50!black},
  fld/.style={->,line width=0.6pt,dashed,cAdB}]

  % --- inputs (column 0) ---
  \node[inp] (eeg)  at (0, 0.00) {EEG $T\!\times\!32$};
  \node[inp] (gaze) at (0,-0.78) {gaze $T\!\times\!6$};
  \node[inp] (imu)  at (0,-1.56) {head $T\!\times\!6$};
  \node[inp] (flow) at (0,-2.34) {flow $T\!\times\!4$};
  \node[inp] (fov)  at (0,-3.12) {fovea $1024$};
  \node[inp] (aud)  at (0,-4.45) {$a_1\dots a_K$};

  % --- encoders (column 1) ---
  \node[enc] (e1) at (3.1, 0.00) {dilated enc};
  \node[enc] (e2) at (3.1,-0.78) {dilated enc};
  \node[enc] (e3) at (3.1,-1.56) {dilated enc};
  \node[enc] (e4) at (3.1,-2.34) {dilated enc};
  \node[frz] (e5) at (3.1,-3.12) {\textbf{frozen} proj};
  \node[enc] (ea) at (3.1,-4.45) {shared audio enc};

  % --- branches (column 2) ---
  \node[nd,fill=cCo,draw=cCoB,minimum width=34mm,minimum height=7.5mm]
       (cpl) at (7.1, 0.00) {$s_k=\tau\,w^{\!\top}\mathrm{corr}_t(e,\hat a_k)$\\[-1pt]
                             \scriptsize centred over $t$ and over $k$};
  \node[nd,fill=cOr,draw=cOrB,minimum width=34mm,minimum height=11mm]
       (ori) at (7.1,-2.10) {speaker heads \emph{(no audio)}\\[-1pt]
                             \scriptsize gate $+$ fusion $\to$ 4 logits};
  \node[cCoB,font=\bfseries\tiny,anchor=south] at (7.1,0.55)
       {COUPLING BRANCH --- \emph{what} is attended};
  \node[cOrB,font=\bfseries\tiny,anchor=north] at (7.1,-2.72)
       {ORIENTATION BRANCH --- \emph{where} attention points};

  % --- adversary ---
  \node[nd,fill=cAd,draw=cAdB,minimum width=34mm,minimum height=5.6mm]
       (adv) at (7.1,-4.45) {audio-only adversary (GRL)};

  % --- scores + loss (columns 3, 4) ---
  \node[nd,fill=gray!10,draw=gray!55,minimum width=17mm,minimum height=9mm]
       (sc) at (10.9,-1.05) {$K$ slot\\scores};
  \node[nd,fill=white,draw=cCoB,minimum width=30mm,minimum height=14mm]
       (loss) at (14.1,-1.05) {cross-entropy\\$+$ contrastive\\[1pt]\scriptsize $+$ EEG null hinge, anti-collapse};

  % --- wiring ---
  \foreach \a/\b in {eeg/e1,gaze/e2,imu/e3,flow/e4,fov/e5,aud/ea} \draw[fl] (\a)--(\b);
  \draw[fl] (e1) -- (cpl);
  \draw[fl] (ea.east) -- ++(0.35,0) |- (cpl.south west);
  \draw[fld] (ea) -- (adv);
  \draw[fl] (e1.east) -- ++(0.18,0) |- (ori.west);
  \foreach \e in {e2,e3,e4,e5} \draw[fl] (\e.east) -- ++(0.18,0) |- (ori.west);
  \draw[fl] (cpl.east) -- ++(0.4,0) |- (sc.north west);
  \draw[fl] (ori.east) -- ++(0.4,0) |- (sc.south west);
  \node[font=\tiny,anchor=west] at (9.05,-1.95) {gather by $\pi$};
  \draw[fl] (sc) -- (loss);
\end{tikzpicture}}
\caption{\textbf{\merit{}.} The coupling branch compares the EEG with each
candidate using a correlation over time with the mean removed. A brain encoder
that outputs the same vector at every step therefore gives every candidate a
score of zero, which is chance. The orientation branch gets \emph{no audio}, so
it cannot use audio cues. The model encodes once and scores separately, so a
trained model can be re-scored under any shuffle of any stream without encoding
again.}
\label{fig:arch}
\end{figure}

\paragraph{Encoders sized to the brain response.}
Each stream goes through five dilated temporal convolutions
\citep{yu2016dilated} with dilations $1,2,4,8,16$ and GroupNorm
\citep{wu2018groupnorm} after each layer. This gives a window of $63$ samples
($0.98$\,s at $64$\,Hz), about the time over which the cortical response to
speech unfolds \citep{lalor2009}. There is no activation on the last layer,
because a rectified output has no negative values and squeezes all the
correlations the score depends on toward one. Padding is symmetric, because the
brain responds $100$--$300$\,ms \emph{after} the sound, so an encoder that looks
only backward in time cannot see the response. Appendix~\ref{app:encoder}
details both choices. One-sided versions restrict the model to a chosen range of
delays (Section~\ref{sec:res-mm}).

\paragraph{A score that cannot use a constant brain signal.}
\label{sec:coupling}
For each of the $D\!=\!16$ embedding dimensions we compute the correlation over
time between the EEG embedding $e$ and the embedding $\hat a_k$ that the shared
audio encoder computes from the $k$-th candidate envelope. We
combine the $D$ values with a weight vector that has no bias, scale by a learned
temperature, and subtract the mean over candidates:
\begin{equation}
  \tilde e=\frac{e-\overline{e}}{\lVert e-\overline{e}\rVert_2},\qquad
  c_d(k)=\textstyle\sum_t \tilde e_{t,d}\,\tilde a_{k,t,d},\qquad
  s_k=\tau\, w^{\!\top} c(k),\qquad
  \mathrm{score}_k=s_k-\tfrac1K\textstyle\sum_j s_j .
  \label{eq:coupling}
\end{equation}
Removing the time average is the key step. If the encoder outputs the same
vector at every time step --- the state it falls into once it learns to ignore
its input --- then $e-\overline e$ is \emph{exactly zero}. Every $c_d(k)$ is
then zero, all $K$ scores are equal, and accuracy is exactly chance. The
shortcut of Section~\ref{sec:failures}(ii) is not discouraged by a penalty; it is
impossible. Two more properties follow. Correlation ignores scale, so a candidate
cannot win by being louder. Subtracting the mean over candidates removes anything
added equally to all of them, which is why $w$ has no bias.

\paragraph{Orientation heads with no audio, and an audio-only adversary.}
Gaze, head motion, scene motion and the foveal stream --- a frozen V-JEPA\,2
embedding of the patch of scene around the gaze point
(Section~\ref{sec:setup}) --- show where the listener is facing. They have no
timed relationship with speech, so asking them to match a speech envelope makes
no sense --- and doing so made them collapse in our first attempt. Instead, each
one feeds a classifier head that predicts the loudspeaker index from the mean
and standard deviation of its embedding over time. These classifier heads get
\textbf{no audio input}, so shuffling them drops accuracy to chance, and anything
above chance is earned by the recording. The EEG also gets such a classifier
head, because attention shifts brain activity toward one side. A learned, non-negative weight
per stream and a small fusion layer combine the loudspeaker scores, which are then
added to the coupling scores. Separately, a small head tries to find the correct
candidate from the audio alone, behind a gradient-reversal layer
\citep{ganin2015}. Any audio cue it can read is thereby trained out of the audio
encoder. This matters because the probe of Section~\ref{sec:failures}(i) is
linear and can miss cues that a network finds.

\paragraph{Training objective.}
Let a batch hold $B$ windows from one listener. For window $i$, $\ell_i\in
\mathbb{R}^K$ are the final candidate scores, $y_i$ is the correct candidate,
$r_i$ is the attended loudspeaker, $e_i$ is the EEG embedding and
$a_i^{+}=\hat a_{i,y_i}$ is the embedding of the correct candidate. We write
$s(e,a)=\tau\,w^{\!\top}c(e,a)$ for the coupling score of Eq.~\ref{eq:coupling}
before centring over candidates, and $\rho$ for a shift of the batch by one
window. \merit{} minimises
\begin{equation}
  \mathcal{L}=\mathcal{L}_{\text{task}}
  +0.3\,\mathcal{L}_{\text{aux}}+\mathcal{L}_{\text{con}}
  +0.5\,\mathcal{L}_{\text{null}}+0.1\,\mathcal{L}_{\text{col}}
  +0.3\,\mathcal{L}_{\text{adv}},
  \label{eq:loss}
\end{equation}
with the six terms
\begin{align*}
  \mathcal{L}_{\text{task}} &= \tfrac1B\textstyle\sum_i
      \mathrm{CE}_{0.1}\big(\ell_i,\,y_i\big),
  \qquad
  \mathcal{L}_{\text{aux}} = \tfrac{1}{B|\mathcal{O}|}\textstyle\sum_i\sum_{m\in\mathcal{O}}
      \mathrm{CE}\big(u_i^m,\,r_i\big),\\
  \mathcal{L}_{\text{con}} &= -\tfrac{1}{2B}\textstyle\sum_i\Big[
      \log\frac{e^{S_{ii}}}{\sum_{j}e^{S_{ij}}}
      +\log\frac{e^{S_{ii}}}{\sum_{j}e^{S_{ji}}}\Big],
  \qquad S_{ij}=s\big(e_i,a_j^{+}\big),\\
  \mathcal{L}_{\text{null}} &= \tfrac1B\textstyle\sum_i\Big[
      \mathrm{sp}\big(s(e_{\rho(i)},a_i^{+})-s(e_i,a_i^{+})+\mu\big)
      +\mathrm{sp}\big(s(\mathbf{0},a_i^{+})-s(e_i,a_i^{+})+\mu\big)\Big],\\
  \mathcal{L}_{\text{col}} &= \tfrac{1}{BD}\textstyle\sum_{i,d}
      \max\big(0,\;\gamma-\mathrm{std}_t\,e_{i,t,d}\big)
      +\tfrac{0.04}{D}\textstyle\sum_{d\neq d'}C_{dd'}^2,\\
  \mathcal{L}_{\text{adv}} &= \tfrac1B\textstyle\sum_i
      \mathrm{CE}\big(\big[g(R(f_{i,k}))\big]_{k=1}^{K},\,y_i\big),
  \qquad f_{i,k}=\big[\,\mathrm{mean}_t\,\hat a_{i,k}\,;\,\mathrm{std}_t\,\hat a_{i,k}\big].
\end{align*}
Here $\mathrm{CE}_{0.1}$ is cross-entropy with label smoothing $0.1$ and
$\mathrm{sp}$ is softplus. $\mathcal{L}_{\text{task}}$ is the decision loss:
the cross-entropy between the $K$ final candidate scores $\ell_i$ and the index
$y_i$ of the attended candidate, so it trains everything that feeds those
scores.
$\mathcal{L}_{\text{aux}}$ gives each orientation head $m\in\mathcal{O}$ its own
loudspeaker prediction $u_i^m$, so the strongest stream cannot take all the
gradient. $\mathcal{L}_{\text{con}}$ is a two-way contrastive loss
\citep{oord2018,radford2021clip} that asks which window's EEG goes with which
window's speech; the batch comes from one listener, so the EEG cannot be matched
by listener identity, and a constant encoder cannot lower this loss.
$\mathcal{L}_{\text{null}}$ requires the real EEG to give the correct candidate a
higher coupling score than another window's EEG or a zeroed EEG does, by a margin
$\mu=0.5$. $\mathcal{L}_{\text{col}}$ keeps the spread over time of each embedding
dimension above $\gamma=0.5$ --- the quantity that falls to zero when an encoder
stops using its input --- and discourages dimensions from copying each other,
where $C$ is the covariance of the time-centred embeddings over all windows and
time steps \citep{bardes2022vicreg}. $\mathcal{L}_{\text{adv}}$ trains a small
network $g$ to find the correct candidate from summary statistics of the audio
embeddings alone; $R$ is a gradient-reversal layer \citep{ganin2015}, which
leaves the forward pass unchanged and flips the sign of the gradient that reaches
the audio encoder, so any audio cue $g$ can read is trained out of it.

\paragraph{Tested but not used: a contribution hinge.}
\label{sec:hinges}
The EEG hinge above acts on the coupling score only. A natural next step is to
apply the same test to the \emph{final} decision, for every stream. For each
stream $m$ we tried
\begin{equation}
  \mathcal{L}^{\text{contrib}}_m=
  \mathrm{softplus}\big(\sigma_y(z_{m\leftarrow\rho})-\sigma_y(z)+\delta\big)
  +\mathrm{softplus}\big(\sigma_y(z_{m\leftarrow 0})-\sigma_y(z)+\delta\big),
  \label{eq:hinge}
\end{equation}
where $\sigma_y(\ell)=\ell_y-\frac1K\sum_k\ell_k$ is the score of the correct
candidate minus the mean score, $z_{m\leftarrow\rho}$ replaces stream $m$ with
another window's copy, $z_{m\leftarrow 0}$ replaces it with zeros, and
$\delta\!=\!0.3$ is the margin. The hinge arms add $0.5\sum_m\lambda_m\mathcal{L}^{\text{contrib}}_m$
to Eq.~\ref{eq:loss}. After each epoch the weight of stream $m$ is set from its
validation contribution $c_m$ in the fused model and the contribution
$c_m^{\text{solo}}$ it reaches through its own head alone:
\begin{equation}
  \lambda_m=\mathrm{clip}\Big(\frac{\max(0,\;c_m^{\text{solo}}-c_m)}{c_m^{\text{solo}}},\;0,\;1\Big),
  \label{eq:gate}
\end{equation}
so a stream that can do well alone but is ignored by the fused model is pushed
hardest. The hinge does not raise the brain's share of credit
(Section~\ref{sec:res-hinge}), so \merit{} trains without it; we report it
because the reason it fails is informative.

\paragraph{Choosing checkpoints by contribution.}
\label{sec:selection}
If a shortcut exists, picking the epoch with the best validation accuracy picks
the epoch that uses the shortcut best. We keep instead the epoch that maximises
\begin{equation}
  \big(v(M)-v(\emptyset)\big) + \beta\,\min_m \phi_m,\qquad \beta=1,
  \label{eq:criterion}
\end{equation}
the validation contribution plus the credit of the \emph{weakest} stream. An epoch that is accurate because it rides one stream loses
to a slightly less accurate epoch that uses all of them. We compute both criteria
at every epoch and report both.

\paragraph{Credit: a split that adds up exactly.}
\label{sec:shapley}
The contributions of single streams do not add up to the total, because streams
overlap --- gaze and the foveal embedding both describe where the listener
looks --- and interact. The Shapley value \citep{shapley1953} over subsets of
\emph{intact} streams is the only split that does add up:
\begin{equation}
  \phi_m=\!\!\sum_{T\subseteq M\setminus\{m\}}\!\!
  \frac{|T|!\,(|M|-|T|-1)!}{|M|!}\Big[v(T\cup\{m\})-v(T)\Big],
  \qquad \sum_{m}\phi_m=v(M)-v(\emptyset).
  \label{eq:shapley}
\end{equation}
This is permutation feature importance \citep{breiman2001,fisher2019,covert2020}
applied to whole streams. With $M\!=\!5$ streams there are only $32$ subsets, so
we compute it exactly, and no subset needs the inputs to be encoded again. The
same random shuffles are used for every subset, so differences between subsets
are not shuffle noise.
